\statement{A.83}{Proved}{Actual odd finite-head twisted duality fixes all self-dual multiplicities}
Let E be the pure negative pants extension in A.82, and
let K be an open normal subgroup of E contained in F.
Put \(Q=E/K\), let \(k=\mathbb F_p\) for an odd
prime p, and define the actual finite module
\[
             \overline R=k\otimes_{\mathbb Z_p}(R_p)_K.
 \tag{A.83.1}
\]
In the Grothendieck group of finite \(k[Q]\)-modules,
\[
                    [\overline R^*]=[\overline R\otimes\varepsilon].
 \tag{A.83.2}
\]
If \(p\nmid|Q|\), this is an actual module isomorphism.
Under this same prime-to-order hypothesis, A.82 and the
isomorphism force every self-dual simple S to
occur in \(\overline R\) with multiplicity
\[
                         2\dim_{\operatorname{End}_Q(S)}S.
 \tag{A.83.3}
\]
In particular, if every simple \(k[Q]\)-module is
self-dual and \(p\nmid|Q|\), then
\(\overline R\simeq k[Q]^2\).
The nonself-dual paired blocks and the modular isomorphism
problem are not removed by this statement.

\begin{proof}
Write
\[
 A=\bigoplus_{i=1}^3 k[Q/q(C_i)],\qquad M=H_1(K,k).
\]
Restriction of every summand of \(P_p\) to K is free,
since K meets its torsion stabilizers trivially. The
homology sequence of (A.82.1), first reduced modulo p,
gives the actual exact sequence
\[
                 0\longrightarrow M\longrightarrow\overline R
                    \longrightarrow A\longrightarrow k
                    \longrightarrow0.
 \tag{A.83.4}
\]
Thus \([\overline R]=[M]+[A]-1\). The module
\(\overline R\) is projective by A.82, although
semisimplicity will only be used under the stated
prime-to-order hypothesis.

\emph{Retain the actual boundary and its sign.}\par
Let \(D_j\) be the actual dihedral normalizer of the
jth base cusp in E, and write \(\overline D_j=q(D_j)\).
Its rotation subgroup is the image of that cusp; its order
may be any positive integer, including one. The actual
cusp permutation module is
\[
                         C=\bigoplus_{j=1}^3 k[Q/\overline D_j].
 \tag{A.83.5}
\]
These cosets count precisely the positive oriented
boundary circles of the K-cover. A negative element
reverses every boundary orientation. The boundary
submodule of M is therefore the quotient of
\(C\otimes\varepsilon\) by its one signed total-sum
relation. If H is the compact-surface homology quotient,
\[
                         [M]=[H]+\varepsilon([C]-1).
 \tag{A.83.6}
\]
This is a quotient by the diagonal total-sum vector,
not a claimed splitting of an augmentation map; it
remains valid when p divides the number of cusps.

The intersection form on H is nondegenerate and changes
by \(\varepsilon\), giving
\[
                                H^*\simeq H\otimes\varepsilon.
 \tag{A.83.7}
\]
The sign is an actual surface-pair calculation, as in C.6
and A.81, rather than a Nielsen realization assumption.
The boundary of the relative degree-two orientation class
is the total sum of the oriented boundary generators.
Literal exponent \(-1\) on them makes the induced
relative degree negative; naturality of the compact
intersection pairing gives (A.83.7). Positive deck
transformations have degree one. This uses the source
cyclic boundary groups and the actual compact quotient.

\emph{Use the finite dihedral polygon before summing
endpoints.}\par
For \(D=\overline D_j\), let \(C_+,C_-\) be the
images of the two local reflection subgroups. A
barycentrically subdivided polygon has two vertex orbits,
one free edge orbit, and orientation character
\(\varepsilon\). Its cellular sequence is
\[
 0\longrightarrow k(\varepsilon)\longrightarrow k[D]
 \longrightarrow k[D/C_+]\oplus k[D/C_-]
 \longrightarrow k\longrightarrow0.
 \tag{A.83.8}
\]
When the rotation image is trivial, the same sequence is
the two-edge polygon, with both vertex stabilizers equal
to D. Thus no separate nondegeneracy assumption on the
finite cusp image is required. In the representation
Grothendieck group, (A.83.8) implies
\[
 (1-\varepsilon)\bigl([k[D/C_+]]+[k[D/C_-]]\bigr)
                              =2(1-\varepsilon)[k].
 \tag{A.83.9}
\]
Here the regular module is unchanged by the sign twist,
and \((1-\varepsilon)^2=2(1-\varepsilon)\).

Induce (A.83.9) to Q and sum over the three actual
cusps. C.28's incidence triangle says that each global
involution class occurs exactly twice among the six
local endpoints. Consequently
\[
                     2(1-\varepsilon)[A]
                              =2(1-\varepsilon)[C].
 \tag{A.83.10}
\]
The group \(G_0(k[Q])\) is free abelian on its simple
modules, so the integer two can be cancelled. We obtain
\((1-\varepsilon)([A]-[C])=0\).

\emph{The precise duality and multiplicity consequence.}\par
Both A and C are permutation modules and hence self-dual.
Using (A.83.4), (A.83.6), and (A.83.7), we find
\[
 \begin{split}
 [\overline R^*]-[\overline R\otimes\varepsilon]
              &=(1-\varepsilon)([A]-[C])\\
              &=0.
 \end{split}
 \tag{A.83.11}
\]
This proves (A.83.2) for every odd p. If \(p\nmid|Q|\),
Maschke's theorem makes equality of these simple
multiplicities an isomorphism; see
\citet[Theorem~7.4 and Proposition~7.9]{AMilne2025}.
A.82 also reduces to
\[
             \overline R\oplus(\overline R\otimes\varepsilon)
                                      \simeq k[Q]^4.
 \tag{A.83.12}
\]
For self-dual S, (A.83.11) equates its two orientation
multiplicities. Their sum in (A.83.12) is four times
the regular multiplicity. This gives (A.83.3), where
\(\operatorname{End}_Q(S)\) is the finite division
field. Over a nonsplitting field the regular multiplicity
is \(\dim_{\operatorname{End}_Q(S)}S\), not
\(\dim_k S\).

For a nonself-dual pair, the duality can instead exchange
opposite orientation imbalances. No such imbalance has
been constructed in an actual \(\overline R\), but
these identities do not exclude it. When p divides
\(|Q|\), the proof asserts the composition-factor
identity (A.83.2) and does not invoke an additional
Cartan-map theorem or a compact inverse-limit duality.
In particular it supplies no unproved full integral
self-duality, finite tree, or marked generation assertion.
A.84 below supplies the additional Cartan-map and compact
projective-dual arguments, upgrading the modular comparison
without removing the nonselfdual ambiguity.
\end{proof}
